\begin{figure}
    \centering
    \includegraphics[width=\linewidth]{figs/data_representation_figure.pdf}
    \captionof{figure}{\textbf{Geometry and Motion representation.}
    % We define both point maps and scene flows in the world coordinate system.
    For a pixel $p_i$ in frame $\bm{I}_i$,
    $\bm{X}_i$ is its corresponding 3D point.
    As this 3D point moves,
    we use $\bm{X}_i^d$ to represent the moved point and $\bm{V}_{i} = (\Delta x, \Delta y, \Delta z)$ to represent the motion.
    Ideally,
    $\bm{X}_i^d$ should align with a matching point $\bm{X}_{i+1}$ in next frame $\bm{I}_{i+1}$.
    However, their pixel indexes are totally different ($p_i$ vs. $p_{i+1}$) and $p_{i+1}$ might even be out of view due to camera/object motion,
    making it impossible to build one-to-one correspondence between $\bm{X}_i^d$ and $\bm{X}_{i+1}$.
    }%
    \label{fig:representation}
    \vspace{-1em}

\end{figure}